\section{Wick 定理及其证明}

本节按 Gross 第 19 章给出粒子--空穴语言下 Wick 定理的完整证明：先证正规序形式（引理 + 归纳），再过渡到时间编序与真空期望。有限温情形在末节用生成函数简述。

\subsection{粒子、空穴与真空}

取 $H_0$ 的基态为填满至 $\varepsilon_F$ 的费米海 $|\Phi_0\rangle$。定义
\begin{equation}
\begin{aligned}
  a_i&=c_i,\quad a_i^\dagger=c_i^\dagger
  &&(\varepsilon_i>\varepsilon_F),\\
  b_i&=c_i^\dagger,\quad b_i^\dagger=c_i
  &&(\varepsilon_i<\varepsilon_F).
\end{aligned}
\end{equation}
则 $a_i|\Phi_0\rangle=b_i|\Phi_0\rangle=0$，故称 $|\Phi_0\rangle$ 为相对费米海的真空。粒子算符与空穴算符相互反对易。相互作用绘景中
$a_i(t)=a_i e^{-i\varepsilon_i t}$，$b_i(t)=b_i e^{-i|\varepsilon_i-\mu|t}$（记号上常把空穴能量取正）。

场算符可分解为产生部分（下标 $+$）与湮灭部分（下标 $-$）：
\begin{equation}
  \psi=\psi^{(+)}+\psi^{(-)},
  \qquad
  \psi^{(-)}|\Phi_0\rangle=0=(\psi^\dagger)^{(-)}|\Phi_0\rangle.
\end{equation}

\subsection{正规序}

粒子/空穴算符乘积的\textbf{正规序} $N(\cdots)$：把全部湮灭算符移到产生算符右边，并乘以完成该重排所需置换的符号（始终按反对易计）。例如（$\varepsilon_i,\varepsilon_k>\varepsilon_F$，$\varepsilon_j<\varepsilon_F$）
\begin{equation}
  N\bigl[c_i(t_1)c_j(t_2)c_k^\dagger(t_3)\bigr]
  =(+)\,c_j(t_2)\,c_k^\dagger(t_3)\,c_i(t_1).
\end{equation}
要点：
\begin{itemize}
  \item 时间依赖不影响正规序的定义；
  \item 产生算符之间（或湮灭算符之间）的次序可任意改写，因它们彼此反对易，所得算符相同；
  \item 对线性组合按线性扩展：$N(\alpha A+\beta B)=\alpha N(A)+\beta N(B)$。
\end{itemize}
因此，任意含正规序因子的真空期望值为零：
\begin{equation}
  \langle\Phi_0|N(\cdots)|\Phi_0\rangle=0,
\end{equation}
因为最右边总有一个湮灭算符作用在 $|\Phi_0\rangle$ 上（或最左边产生算符作用在 $\langle\Phi_0|$ 上）。

\subsection{配对（pairing）}

对两个粒子/空穴算符（或其线性组合的分量）定义
\begin{equation}
  \overline{AB}
  :=AB-N(AB).
\label{eq:pairing}
\end{equation}
若 $\{A,B\}=0$，则 $N(AB)=-BA$，从而 $\overline{AB}=0$。非零配对必为 $c$-数。显式地，
\begin{equation}
\begin{aligned}
  \overline{a_i(t)a_j^\dagger(t')}
  &=\delta_{ij}\,e^{-i\varepsilon_i(t-t')},\\
  \overline{b_i(t)b_j^\dagger(t')}
  &=\delta_{ij}\,e^{-i|\varepsilon_i|(t-t')},
\end{aligned}
\end{equation}
其余粒子/空穴组合的配对为零。由此
\begin{equation}
  \langle\Phi_0|AB|\Phi_0\rangle
  =\overline{AB}.
\end{equation}

含配对的正规序定义为：先把配成对的算符移到最左（计符号 $(-1)^q$），再对剩余算符取正规序；配对本身是 $c$-数，可提出。

\subsection{引理}

\begin{lemma}[右乘湮灭/产生算符]
对任意粒子/空穴算符 $A_1,\ldots,A_n$ 与 $B$，
\begin{equation}
  N(A_1\cdots A_n)\,B
  =\sum_{r=1}^n
  N\bigl(A_1\cdots\overline{A_r B}\cdots A_n\bigr)
  +N(A_1\cdots A_n B).
\label{eq:wick_lemma}
\end{equation}
\end{lemma}

\begin{proof}
若 $B$ 为湮灭算符，则一切 $\overline{A_r B}=0$，引理化为
$N(A_1\cdots A_n)B=N(A_1\cdots A_n B)$，由正规序定义成立。

若 $B$ 为产生算符，则仅当 $A_r$ 为湮灭算符时 $\overline{A_r B}$ 可能非零。此时只需对“$A_1,\ldots,A_n$ 全为湮灭算符”证明
\begin{equation}
  A_1\cdots A_n B
  =\sum_{r=1}^n
  (-1)^{r+n}\,\overline{A_r B}\,
  A_1\cdots\widehat{A_r}\cdots A_n
  +(-1)^n B A_1\cdots A_n,
\label{eq:lemma_ann}
\end{equation}
（$\widehat{A_r}$ 表示删去 $A_r$；符号来自把 $A_r$ 与 $B$ 移到一起所需的交换。）一般情形可在左边再乘产生算符还原。

对 \eqref{eq:lemma_ann} 用归纳法。$n=1$：由 \eqref{eq:pairing} 即
$A_1 B=\overline{A_1 B}+N(A_1 B)=\overline{A_1 B}-B A_1$。
设 $n$ 成立，左边乘湮灭算符 $A_0$：
\begin{equation}
\begin{aligned}
  A_0 A_1\cdots A_n B
  &=\sum_{r=1}^n(-1)^{r+n}\overline{A_r B}\,
  A_0 A_1\cdots\widehat{A_r}\cdots A_n\\
  &\quad+(-1)^n A_0 B\,A_1\cdots A_n.
\end{aligned}
\end{equation}
再用 $A_0 B=-B A_0+\overline{A_0 B}$，最后一项拆成
$(-1)^n\overline{A_0 B}\,A_1\cdots A_n+(-1)^{n+1}B A_0\cdots A_n$，
恰为 $n+1$ 时的 \eqref{eq:lemma_ann}。
\end{proof}

\subsection{正规序形式的 Wick 定理}

\begin{theorem}[Wick，普通乘积]
\begin{equation}
\begin{aligned}
  A_1\cdots A_n
  &=N(A_1\cdots A_n)\\
  &\quad+\sum_{i<j}
  N\bigl(A_1\cdots\overline{A_i A_j}\cdots A_n\bigr)\\
  &\quad+\sum_{\text{two pairs}}
  N(\cdots\overline{\,\cdot\,}\,\overline{\,\cdot\,}\cdots)
  +\cdots
  +\text{(even $n$) fully paired terms}.
\end{aligned}
\label{eq:wick_N}
\end{equation}
\end{theorem}

\begin{proof}
对 $n$ 归纳。$n=1,2$ 即配对定义。设对 $n$ 成立，右边乘 $A_{n+1}$：
\begin{equation}
  A_1\cdots A_n A_{n+1}
  =N(A_1\cdots A_n)A_{n+1}
  +\sum_{\text{terms with pairings}}N(\cdots)\,A_{n+1}.
\end{equation}
对每一项用引理 \eqref{eq:wick_lemma}：把 $A_{n+1}$ 吞进正规序时，要么进入 $N(\cdots A_{n+1})$，要么与某个 $A_r$ 结成新配对。前者给出 $(n+1)$ 个算符的无配对正规序；后者把“$k$ 个配对”升级为“$k+1$ 个配对”。全部项恰好穷尽 \eqref{eq:wick_N} 在 $n+1$ 时的内容。

由线性性，定理对粒子/空穴算符的线性组合（因而对场算符）同样成立。
\end{proof}

\subsection{收缩与时间编序}

对含时算符定义\textbf{收缩}（contraction）
\begin{equation}
  \overline{A(t)B(t')}
  :=
  T\bigl(A(t)B(t')\bigr)
  -N\bigl(A(t)B(t')\bigr).
\label{eq:contraction}
\end{equation}
由 $T$ 与 $N$ 的反对称性，
\begin{equation}
  \overline{A(t)B(t')}
  =-\,\overline{B(t')A(t)}.
\end{equation}
与配对的关系为：当 $t>t'$ 时收缩等于配对 $\overline{A(t)B(t')}$；当 $t'<t$ 时收缩等于配对 $\overline{B(t')A(t)}$。故收缩仍是 $c$-数，且
\begin{equation}
  \langle\Phi_0|T\bigl(A(t)B(t')\bigr)|\Phi_0\rangle
  =\overline{A(t)B(t')}.
\end{equation}
对 $H_0$ 的产生湮灭算符，非零收缩正是自由传播子：
\begin{equation}
\begin{aligned}
  \overline{c_j(t)c_k^\dagger(t')}
  &=\delta_{jk}\,e^{-i\varepsilon_j(t-t')}
  \bigl[
  \theta(t-t')\theta(\varepsilon_j-\varepsilon_F)
  -\theta(t'-t)\theta(\varepsilon_F-\varepsilon_j)
  \bigr]\\
  &=i G^{(0)}(j t,k t').
\end{aligned}
\label{eq:contr_G0}
\end{equation}
场算符同理：$\overline{\psi(xt)\psi^\dagger(yt')}=iG^{(0)}(xt,yt')$，而
$\overline{\psi\psi}=\overline{\psi^\dagger\psi^\dagger}=0$。

\subsection{时间编序形式的 Wick 定理}

\begin{theorem}[Wick，$T$ 乘积]
\begin{equation}
\begin{aligned}
  T(A_1\cdots A_n)
  &=N(A_1\cdots A_n)\\
  &\quad+\sum_{i<j}
  N\bigl(A_1\cdots\overline{A_i A_j}\cdots A_n\bigr)\\
  &\quad+\cdots
  +\text{fully contracted terms (even $n$)}.
\end{aligned}
\label{eq:wick_T}
\end{equation}
费米子：把一对收缩算符移到乘积最左所需交换次数为奇时，该项带负号。
\end{theorem}

\begin{proof}
设 $P$ 为把时间排成 $t_{P(1)}\ge t_{P(2)}\ge\cdots$ 的置换，则
\begin{equation}
  T(A_1\cdots A_n)
  =\mathrm{sgn}(P)\,
  A_{P(1)}\cdots A_{P(n)}.
\end{equation}
对右边用已证的 \eqref{eq:wick_N}。因时间已排序，其中每个配对恰等于对应的收缩 \eqref{eq:contraction}。再把下标 $P(i)$ 写回 $1,\ldots,n$，并吸收 $\mathrm{sgn}(P)$ 与费米交换符号，即得 \eqref{eq:wick_T}。
\end{proof}

\subsection{真空期望：完全收缩}

对 \eqref{eq:wick_T} 取 $\langle\Phi_0|\cdots|\Phi_0\rangle$，一切含 $N(\cdots)$ 的项消失，故
\begin{equation}
  \langle\Phi_0|T(A_1\cdots A_{2m})|\Phi_0\rangle
  =\sum_{\mathrm{perfect\,matchings}}
  (-1)^P
  \prod_{\mathrm{pairs}}
  \overline{A_{i_s}A_{j_s}}.
\label{eq:wick_vev}
\end{equation}
奇数个算符的真空期望为零。这就是图技术的出发点：微扰展开中的
$\langle T\,V_I(t_1)\cdots\psi\psi^\dagger\rangle_0$ 化为 $G^{(0)}$ 线的产物。

\subsection{四点函数示例}

\begin{equation}
\begin{aligned}
  &\langle\Phi_0|
  T\,c_1 c_2 c_3^\dagger c_4^\dagger
  |\Phi_0\rangle\\
  &\quad=
  \overline{c_1 c_4^\dagger}\,\overline{c_2 c_3^\dagger}
  -\overline{c_1 c_3^\dagger}\,\overline{c_2 c_4^\dagger}.
\end{aligned}
\end{equation}
第一项为直接收缩，第二项为交换收缩（多一次费米交换）。用 \eqref{eq:contr_G0} 即得两点 $G^{(0)}$ 的行列式结构。

\subsection{有限温度：生成函数证明概要}

巨正则自由理论中，源耦合的生成函数是高斯的（Mahan）：
\begin{equation}
  \Bigl\langle
  T_\tau\exp\Bigl(
  \sum_\alpha
  \bigl[
  \eta_\alpha^*(\tau)c_\alpha(\tau)
  +c_\alpha^\dagger(\tau)\eta_\alpha(\tau)
  \bigr]
  \Bigr)
  \Bigr\rangle_0
  =\exp\bigl(\eta^*\,\mathcal{G}^{(0)}\,\eta\bigr).
\end{equation}
对 Grassmann 源 $\eta,\eta^*$ 求 $2m$ 次导并令源为零，左边给出
$\langle T_\tau c\cdots c^\dagger\rangle_0$，右边给出全部两两配对 $\mathcal{G}^{(0)}$ 的求和（含费米符号）。这与 \eqref{eq:wick_vev} 同构，只是真空换为自由热平均、收缩换为 $-\mathcal{G}^{(0)}$。

\begin{note}
Wick 定理依赖自由（二次型）动力学：只有此时配对/收缩为 $c$-数且高阶矩由二阶矩决定。有相互作用的热力学平均一般不能直接 Wick 分解；相互作用绘景微扰展开中，对 $H_0$ 演化下的编序平均，本定理仍然适用。
\end{note}
